%% ============================================================
%%  Section: Related Work
%%  Skeleton — \placeholder{} marks references and summaries needed.
%% ============================================================

\section{Related Work}
\label{sec:related_work}

\paragraph{Generalization in Visual RL.}
Generalization in vision-based reinforcement learning has been studied extensively
under procedural environment generation~\citep{procgen2020} and visual perturbations.
Data augmentations such as DrQ-v2~\citep{drqv2} and SVEA~\citep{svea2021} improve empirical
robustness, yet predicting policy failure without environment interactions remains open.

\paragraph{Out-of-Distribution Detection and Representation Drift.}
Classical out-of-distribution (OOD) detection in supervised vision relies on density
estimation or feature distance~\citep{\placeholder{cite_ood_detection}}. However, in visual RL,
feature shifts do not uniformly induce reward degradation; background distractors may create
substantial feature drift while leaving task competence unaffected.

\paragraph{Test-Time Training and Adaptation.}
Test-Time Training (TTT)~\citep{sun2020ttt} and TTT++~\citep{liu2021tttplusplus} adapt model
representations at test time using self-supervised objectives such as contrastive or moment
matching losses. In control domains, unconstrained adaptation risks catastrophic representation
collapse, motivating safeguards and trust-region validation~\citep{\placeholder{cite_safeguarded_ttt}}.
